\section{Root Tag Hash Security and Committing Security}
\label{sec:cmtds3}
\label{sec:root-tag-hash-security}

This section proves the digest claim for the root of the \TW{} tree. The
construction adapts the KangarooTwelve final-node pattern to authenticated
encryption: leaf tags play the role of chaining values, and the root emits the
single tag returned by $\HTW$. The proof therefore tracks the final root inputs
that can matter in the hash and committing games. The Root Tag Candidate
Sequence lemma (\Cref{lem:root-candidate-seq}) collects those inputs into a
chronological candidate sequence, and the Root Tag Hash and Target-Hit Bound
(\Cref{lem:root-tag-hash}) turns a collision or target hit among the
candidates into the matching adaptive bound of \Cref{app:lemmas}, controlling
the chance that distinct inputs land on the same $\tauroot$-bit root tag
projection.

The hash random oracle theorems---collision (\Cref{thm:coll-ro}) and
preimage (\Cref{thm:pre-ro,thm:epre-ro})---are then short instantiations,
differing only in which case of the bound they invoke, whether the target is
external or sampled from a fixed input slice, and whether a leaf tag collision
term is needed when distinct inputs may use different ciphertext bodies. CMTDs-4
(\Cref{thm:cmtds4-ro}) and the remaining committing notions reuse the collision
case.
\Cref{app:flat-sponge-lift} transfers each instance to the public
permutation setting, \Cref{sec:concrete-security} records the concrete bounds,
and \Cref{cor:committing-combined} forwards the remaining \cite{BH22}
committing notions to the same root-collision bound.

The instantiations rest on two structural separation facts, recorded
first: distinct contexts $(K, U, A)$ reach distinct final root
transcripts regardless of the leaf tags they aggregate (Opening Triple
Separation, \Cref{cor:triple-sep}), and---absent a collision among those
leaf tags---so do distinct full inputs (Final-Root Input Separation,
\Cref{prop:final-root-sep}).

\subsection{Structural Separation}

\begin{proposition}[Opening Triple Separation]
\label{cor:triple-sep}
If two \TW{} root computations have distinct $(K, U, A)$ triples, then
their bridged final root transcript inputs under $\flatmap(\cdot)$ are
distinct, regardless of whether each computation is an encryption or
a decryption check, and whether or not they act on a common ciphertext.
\end{proposition}

\begin{proof}
Suppose for contradiction that $\flatmap(R_i) = \flatmap(R_j)$. Since the
bridge $\flatmap(\cdot)$ is injective on well-formed transcript prefixes
(\Cref{lem:bridge-decode}), the native flattened inputs of the two
final-root prefixes are equal. Then \Cref{lem:transcript-inj}
recovers the full duplex-call sequence and reads off the pair
$(\sigma,\sigph)$ for each call. The seven $4$-bit phase suffixes are
pairwise distinct, so the \siginit, $\sigma_{\mathsf{ad}}^\pm$,
$\sigma_{\mathsf{msg}}^\pm$, and (when present) $\sigma_{\mathsf{agg}}^\pm$
calls in the recovered sequence are unambiguously identified. The
aggregation phase is present exactly when $n \geq 1$, and the two
recovered sequences are equal, so they either both carry an aggregation
phase or both omit it; either way the recovery of $(K, U, A)$ below uses only the
\siginit and AD calls, which precede any message or aggregation block.

Equality of the recovered sequences forces equality of the \siginit
call, whose $\sigma$ is $\prfxroot \concat K \concat U$ with
fixed-width $K$ and $U$, hence $K_i = K_j$ and $U_i = U_j$. It also
forces equality of the recovered AD subsequences. By the AD bullet of
\Cref{lem:transcript-inj} the associated data phase contributes a
$\sigadlast$-tagged block exactly when the associated data is nonempty,
so the recovered sequences agree on the presence of the AD phase. If
exactly one of $A_i, A_j$ were empty the recovered sequences would
differ in the presence of a $\sigadlast$ call, contradicting their
equality; hence either both are empty, giving $A_i = A_j = \varepsilon$,
or both are nonempty, in which case the phase-recovery sub-claim of
\Cref{lem:transcript-inj} inverts the $\rho$-byte-block partition map on
the AD byte string and equal AD $\sigma$-block sequences imply
$A_i = A_j$. In all cases $(K_i, U_i, A_i) = (K_j, U_j, A_j)$,
contradicting the hypothesis of distinct triples.
\end{proof}

Triple Separation handles distinct contexts; when inputs share a context
but differ in their messages, distinctness of the final root transcripts
needs the absence of a leaf tag collision, recorded next.

\begin{proposition}[Final-Root Input Separation]
\label{prop:final-root-sep}
Let $\badleaf^\star$ be the event that two distinct construction-generated
final leaf transcripts receive the same $\tauleaf$-bit $\RF$ projection.
On $\neg\badleaf^\star$, if two \TW{} construction computations reach
final root transcripts with $\flatmap(R) = \flatmap(R')$, then they share
the same context $(K, U, A)$ and the same full ciphertext $C$.
Equivalently, on $\neg\badleaf^\star$ computations with distinct
$(K, U, A, C)$ reach distinct bridged final root transcripts.
\end{proposition}

\begin{proof}
By bridge injectivity (\Cref{lem:bridge-decode}),
$\flatmap(R) = \flatmap(R')$ gives equal native flattened final-root
inputs. \Cref{lem:transcript-inj} then recovers the same root
initialization $\prfxroot \concat K \concat U$---hence the same $K$ and
$U$---the same associated data blocks---hence the same $A$---the same root
ciphertext chunk, and the same leaf count and ordered leaf tag sequence.
Under $\neg\badleaf^\star$, Leaf-Transcript Recovery
(\Cref{lem:leaf-recovery}) identifies the equal leaf tag sequences with
the same final leaf transcripts, hence the same leaf ciphertext bodies;
with the common root ciphertext chunk this gives the same full $C$.
\end{proof}

\subsection{Root Tag Candidate Bound}

\begin{lemma}[Root Tag Candidate Sequence]
\label{lem:root-candidate-seq}
Consider a random oracle game of \Cref{sec:ro-model} in which the adversary
$\mathcal{B}$ makes direct $\ROflat$ queries and the challenger runs
$m \geq 1$ deterministic encryption or decryption checks at fixed points in
the game, the $i$-th reaching a final root transcript $R_i$ whose root tag is
read by a single $\ROflat$ lookup at $\flatmap(R_i)$. Every execution then
admits a chronological sequence of root tag inputs
such that
\begin{enumerate}
\item each distinct direct query, when it is made, whose input $Y$ decodes under
  $\KeyedTWOut$ to a root tag output-point class---$\Rmsg^+$ (the final
  root message block, which carries the root tag only on the no-leaf path)
  or $\Rtag$ (the aggregated root tag)---is appended as a candidate before
  its $\ROflat$ lookup, and contributes at most one candidate; multi-chunk
  computations also reach $\Rmsg^+$ with zero requested output before
  aggregation, so this predicate over-approximates the root tag candidate
  set, which only loosens the count;
\item if the challenger reaches its checks, each $\flatmap(R_i)$ is in the
  sequence, appended---if not already present---immediately before the
  $i$-th root tag answer is sampled;
\item no candidate's $\ROflat$ value is read before its append step, so
  the sequence satisfies the predictability and unrevealed-value premises
  of \Cref{lem:adaptive-candidate};
  and
\item its length is at most $\qRO(\mathcal{B}) + m$.
\end{enumerate}
\end{lemma}

\begin{proof}
The committing and hash games have no persistent challenger-sampled key: all
keys are adversary-supplied except for the hidden source in the standard
preimage game. A direct query to a keyed transcript is therefore one of
$\mathcal{B}$'s visible $\ROflat$ queries counted in $\qRO(\mathcal{B})$,
not a separate construction query. Each direct-query candidate is appended
before that direct query reads its value. The append predicate depends on $Y$
alone through the exact decoder of \Cref{lem:bridge-decode}, which ignores the
requested output length, and by Output-Point Separation
(\Cref{cor:output-sep}) each input decodes to at most one output-point class;
hence each distinct direct query contributes at most one candidate,
giving~(1).

If the challenger reaches its $m$ checks, append $\flatmap(R_i)$
immediately before the $i$-th root tag answer is sampled, unless it is
already present. The intermediate $\ROflat$ lookups of check $i$ lie at
non-root tag output points, so by \Cref{cor:output-sep} they do not read
$\flatmap(R_i)$. If $\flatmap(R_i)$ was touched earlier---by a direct
query or by an earlier check $j < i$---then \Cref{lem:bridge-decode} and
\Cref{cor:output-sep} make that earlier access a final-root lookup at the
same address, which was therefore already appended, so the input is in
the sequence and is not appended again. Otherwise its value is still
unread when appended. Either way no candidate's value is read before its
append step, giving~(2) and~(3). At most $\qRO(\mathcal{B})$ direct
candidates and $m$ challenger candidates are appended, giving~(4).
\end{proof}

\begin{lemma}[Root Tag Hash and Target-Hit Bound]
\label{lem:root-tag-hash}
Consider a \TW{} random oracle game (\Cref{sec:ro-model}) with no
persistent challenger-sampled key---all keys adversary-supplied except,
in the standard preimage game, the hidden source
(\Cref{lem:root-candidate-seq})---in
which $\mathcal{B}$ makes direct $\ROflat$ queries and the challenger runs
deterministic encryption or decryption checks reaching final root transcripts,
with win event $\mathsf{W}$. By the Root Tag Candidate Sequence lemma
(\Cref{lem:root-candidate-seq}) the root tag inputs form a chronological
candidate sequence meeting the premises of \Cref{lem:adaptive-candidate} with
$\tau = \tauroot$. Writing
$\qRO = \qRO(\mathcal{B})$:
\begin{enumerate}
\item if $\mathsf{W}$ forces $s$ checks to reach pairwise distinct bridged
  final root transcripts carrying a common root tag, then, with $m = s$
  checks, $\Pr[\mathsf{W}] \leq \RootColl_{\tauroot}(\qRO, s)$;
\item if $\mathsf{W}$ forces a single check ($m = 1$) to carry a root tag
  equal to a target fixed independently of $\ROflat$, then
  $\Pr[\mathsf{W}] \leq \RootPre_{\tauroot}(\qRO)$;
\item if the game first runs a designated source check, reveals its root tag,
  and $\mathsf{W}$ forces one later check to reach a distinct bridged final
  root transcript carrying that designated tag, then, with $m = 2$ checks,
  $\Pr[\mathsf{W}] \leq \RootPre_{\tauroot}(\qRO)$.
\end{enumerate}
\end{lemma}

\begin{proof}
In each case the candidate sequence has the stated length and meets the
premises of \Cref{lem:adaptive-candidate}.
\emph{(1)} The $s$ distinct candidates carrying a common
$\tauroot$-projection are an $s$-way collision; the collision case with
$L = \qRO + s$ gives
$\binom{\qRO + s}{s} \cdot 2^{-\tauroot(s-1)}
= \RootColl_{\tauroot}(\qRO, s)$.
\emph{(2)} The candidate hitting the fixed target is a preimage; the
preimage case with $L = \qRO + 1$ gives
$(\qRO + 1) \cdot 2^{-\tauroot} = \RootPre_{\tauroot}(\qRO)$.
\emph{(3)} Let $j_0$ be the position of the designated source check. The
later check is required to be a distinct bridged final root transcript, so it
is a candidate position $j \neq j_0$ with the same $\tauroot$-bit projection.
The candidate sequence has length at most $L = \qRO + 2$, and the
second-preimage corollary of \Cref{lem:adaptive-candidate} gives
$(L-1)2^{-\tauroot} = (\qRO+1)2^{-\tauroot}
= \RootPre_{\tauroot}(\qRO)$.
\end{proof}

\subsection{Hash Random Oracle Bounds}

\begin{theorem}[Random Oracle Collision Resistance of $\HTW$]
\label{thm:coll-ro}
For every $s \geq 2$ and every $s$-way collision
adversary $\mathcal{B}$ in the \TW{} random oracle model,
\[
  \AdvColl_{\RO}(\mathcal{B})
  \;\leq\;
  \LeafColl_{\tauleaf}(\qRO(\mathcal{B})+\nleaf(\mathcal{B}))
  + \RootColl_{\tauroot}(\qRO(\mathcal{B}), s).
\]
\end{theorem}

\begin{proof}
Run the random oracle collision game. Let $\badleaf^\star$ be the event
that two distinct construction-generated final leaf transcripts receive
the same $\tauleaf$-bit $\RF$ projection; by the general count
of \Cref{lem:leaf-collision},
$\Pr[\badleaf^\star] \leq \LeafColl_{\tauleaf}(\qRO(\mathcal{B})+\nleaf(\mathcal{B}))$,
so it remains to bound the win under $\neg\badleaf^\star$. After rejecting
unless the $s$ output inputs are pairwise distinct and in the \TW{}
domain, the challenger performs $s$ deterministic encryptions
$\Enc_{K_i}^{\RO}(U_i,A_i,M_i)$, reaching $R_1,\dots,R_s$, all carrying a
common root tag $T$ on a win. Pairwise distinct tuples $(K_i,U_i,A_i,M_i)$
have pairwise distinct $(K_i,U_i,A_i,C_i)$: two tuples differing in their
context already differ in $(K,U,A)$, while two sharing a context but
differing in the message yield distinct ciphertexts, since for a fixed
context encryption maps the message to the ciphertext bijectively. So on
$\neg\badleaf^\star$ Final-Root Input Separation
(\Cref{prop:final-root-sep}) makes $\flatmap(R_1),\dots,\flatmap(R_s)$
pairwise distinct. The Root Tag Hash and Target-Hit Bound (\Cref{lem:root-tag-hash},
part~1) bounds the win under $\neg\badleaf^\star$ by
$\RootColl_{\tauroot}(\qRO(\mathcal{B}),s)$; adding $\Pr[\badleaf^\star]$
gives the theorem.
\end{proof}

\begin{theorem}[Random Oracle Preimage Resistance of $\HTW$]
\label{thm:pre-ro}
Fix byte lengths $0 \leq a \leq A_{\max}$ and
$0 \leq \mu \leq M_{\max}$, and let
$m_{a,\mu}=k+\nu+8a+8\mu$. For every Pre adversary
$\mathcal{B}$ against $\HTW$ on the fixed input slice
$\mathcal{X}_{a,\mu}$ in the \TW{} random oracle model
(\Cref{sec:ro-model}),
\[
\begin{aligned}
  \AdvPreSlice{a}{\mu}_\RO(\mathcal{B})
  \;\leq\;&
  \LeafColl_{\tauleaf}(\qRO(\mathcal{B})+\nleaf(\mathcal{B})) \\
  &+ \RootPre_{\tauroot}(\qRO(\mathcal{B}))
  + 2^{\tauroot-m_{a,\mu}} .
\end{aligned}
\]
\end{theorem}

\begin{proof}
Run the random oracle Pre game. The challenger samples
$X^* \sample \mathcal{X}_{a,\mu}$, writes
$X^*=(K^*,U^*,A^*,M^*)$, computes $T^*=\HTW(X^*)$, gives $T^*$ to
$\mathcal{B}$, and later checks the output $X=(K,U,A,M)$. Let
$\mathsf{same}$ be the event $X=X^*$. For a
fixed random oracle table and fixed adversary coins, the adversary's output
is a function of the $\tauroot$-bit challenge tag $T^*$, so over the
uniform source $X^*$ it can equal $X^*$ with probability at most
$2^{\tauroot}/2^{m_{a,\mu}}=2^{\tauroot-m_{a,\mu}}$. Hence
$\Pr[\mathsf{same}] \leq 2^{\tauroot-m_{a,\mu}}$.

It remains to bound winning executions with $X \neq X^*$. Let
$\badleaf^\star$ be the event that two distinct construction-generated
final leaf transcripts in the source and output checks receive the same
$\tauleaf$-bit $\RF$ projection. By the general count of
\Cref{lem:leaf-collision}---which charges every direct query and every
construction-generated final leaf transcript regardless of key
provenance, covering the challenger-sampled source key---
\[
  \Pr[\badleaf^\star]
  \leq
  \LeafColl_{\tauleaf}(\qRO(\mathcal{B})+\nleaf(\mathcal{B})).
\]
On $\neg\badleaf^\star$, the distinct tuples $X^*$ and $X$ have distinct
final root transcripts: distinct contexts are separated by
\Cref{cor:triple-sep}, while equal contexts and distinct messages produce
distinct ciphertexts, and then \Cref{prop:final-root-sep} applies. A win
with $X \neq X^*$ therefore forces the later output check to hit the
designated source tag $T^*$ at a distinct bridged final root transcript.
The sampled-source case of \Cref{lem:root-tag-hash} bounds this probability
by $\RootPre_{\tauroot}(\qRO(\mathcal{B}))$. Adding the source-guess and
leaf-collision events gives the theorem.
\end{proof}

\begin{theorem}[Random Oracle Everywhere Preimage Resistance of $\HTW$]
\label{thm:epre-ro}
For every query-independent target tag $T^* \in \bits^{\tauroot}$ and
every ePre adversary $\mathcal{B}$ in the \TW{}
random oracle model (\Cref{sec:ro-model}),
\[
  \AdvePre_\RO(\mathcal{B})
  \;\leq\; \RootPre_{\tauroot}(\qRO(\mathcal{B}))
  = \frac{\qRO(\mathcal{B}) + 1}{2^{\tauroot}}.
\]
\end{theorem}

\begin{proof}
Run the random oracle preimage game for the query-independent target
$T^*$, which is independent of $\ROflat$. After rejecting unless the
output input $X = (K,U,A,M)$ is in the \TW{} domain, the challenger
performs one encryption $\Enc_{K}^\RO(U,A,M)$, reaching a final root
transcript $R$. A win means $\HTW(X) = T^*$, i.e.\ the root tag
$\lfloor \ROflat(\flatmap(R)) \rfloor_{\tauroot}$ at $R$ equals the fixed
target $T^*$. The Root Tag Hash and Target-Hit Bound (\Cref{lem:root-tag-hash}, part~2,
with target $T^*$) gives
$\AdvePre_\RO(\mathcal{B}) \leq \RootPre_{\tauroot}(\qRO(\mathcal{B}))$.
\end{proof}

\subsection{Public Permutation Hash Bounds}
\label{sec:root-tag-pp}

The flat sponge lift transfers the preceding random oracle bounds to the public
permutation model. For these one-world hash games, the lift contributes
$\Lift_c(\Ntot)$, and the derived random oracle adversary has
$\qRO(\mathcal{B}) \leq \Nprim$; challenger construction
computations---including the preimage game's source encryption---are
construction-internal work counted in $N$.

\begin{theorem}[Collision Resistance of $\HTW$]
\label{thm:coll}
For every $s \geq 2$ and every adversary $\mathcal{A}$ against the $s$-way
multi-collision resistance of $\HTW$ in the sense of \cite{BH22}, with the
resources of \Cref{sec:resources},
\[
  \AdvColl_\TW(\mathcal{A})
  \;\leq\;
  \Lift_c(\Ntot)
  + \LeafColl_{\tauleaf}(\Nprim+\Nleaf)
  + \RootColl_{\tauroot}(\Nprim, s).
\]
\end{theorem}

\begin{proof}
By the Lift Application corollary (\Cref{cor:lift-application}) for the
collision game, the derived random oracle adversary $\mathcal{B}$ has
$\qRO(\mathcal{B}) \leq \Nprim$ and $\nleaf(\mathcal{B}) \leq \Nleaf$. The
random oracle bound of \Cref{thm:coll-ro} is nondecreasing in both
$\qRO(\mathcal{B})$ and $\nleaf(\mathcal{B})$, so the monotone form of
\Cref{cor:lift-application} gives the stated bound. The $s$ colliding
inputs may differ in their messages, hence in their ciphertext bodies; the
$\LeafColl$ term accounts for the only way two distinct inputs sharing a
context $(K, U, A)$ can reach the same final root transcript---a collision
among the leaf tags it absorbs---while distinct contexts already reach
distinct final root transcripts by Opening Triple Separation
(\Cref{cor:triple-sep}).
\end{proof}

\begin{corollary}[Two-Way Collision Resistance of $\HTW$]
\label{cor:coll-rs04}
Specializing \Cref{thm:coll} to $s = 2$ gives the classical two-way
collision resistance of \cite{RS04}: for every collision adversary
$\mathcal{A}$ against $\HTW$ with the resources of \Cref{sec:resources},
\[
  \AdvColl_\TW(\mathcal{A})
  \;\leq\;
  \Lift_c(\Ntot)
  + \LeafColl_{\tauleaf}(\Nprim+\Nleaf)
  + \RootColl_{\tauroot}(\Nprim, 2),
\]
with $\RootColl_{\tauroot}(\Nprim, 2) = \binom{\Nprim + 2}{2} / 2^{\tauroot}$.
\end{corollary}

\begin{proof}
The $s = 2$ case of \Cref{thm:coll}; \cite{BH22} note that $s$-way
multi-collision resistance at $s = 2$ is the classical collision resistance
of \cite{RS04}.
\end{proof}

\begin{corollary}[Second-Preimage Resistance of $\HTW$]
\label{cor:second-preimage-rs04}
For every adversary against $\HTW$, with the resources of
\Cref{sec:resources}, in either the standard second-preimage (Sec) or
everywhere second-preimage (eSec) sense of \cite{RS04}, the corresponding
advantage is bounded by the two-way collision bound of
\Cref{cor:coll-rs04}:
\[
  \Lift_c(\Ntot)
  + \LeafColl_{\tauleaf}(\Nprim+\Nleaf)
  + \RootColl_{\tauroot}(\Nprim, 2).
\]
\end{corollary}

\begin{proof}
By \cite[Proposition~6]{RS04}, collision resistance implies both Sec and
eSec. Applying this implication to the two-way collision bound of
\Cref{cor:coll-rs04} gives the stated bound.
\end{proof}

\begin{theorem}[Preimage Resistance of $\HTW$]
\label{thm:pre}
Fix byte lengths $0 \leq a \leq A_{\max}$ and
$0 \leq \mu \leq M_{\max}$, and let
$m_{a,\mu}=k+\nu+8a+8\mu$. For every adversary against $\HTW$, with the
resources of \Cref{sec:resources}, in the standard preimage (Pre) sense of
\cite{RS04} on the fixed input slice $\mathcal{X}_{a,\mu}$,
\[
\begin{aligned}
  \AdvPreSlice{a}{\mu}_\TW(\mathcal{A})
  \;\leq\;&
  \Lift_c(\Ntot)
  + \LeafColl_{\tauleaf}(\Nprim+\Nleaf) \\
  &+ \RootPre_{\tauroot}(\Nprim)
  + 2^{\tauroot - m_{a,\mu}}.
\end{aligned}
\]
\end{theorem}

\begin{proof}
By the Lift Application corollary (\Cref{cor:lift-application}) for the
preimage game, the derived random oracle adversary $\mathcal{B}$ has
$\qRO(\mathcal{B}) \leq \Nprim$ and
$\nleaf(\mathcal{B}) \leq \Nleaf$; per \Cref{sec:resources}, these caps
cover both challenger encryptions, the source encryption's duplex calls
and final leaf transcripts included. The random oracle bound of
\Cref{thm:pre-ro} is nondecreasing in both resources, so the monotone form
of \Cref{cor:lift-application} gives the stated bound.
\end{proof}

\begin{theorem}[Everywhere Preimage Resistance of $\HTW$]
\label{thm:epre}
For every target tag $T^* \in \bits^{\tauroot}$ fixed in advance and every
ePre adversary $\mathcal{A}$ against $\HTW$ with the resources of
\Cref{sec:resources},
\[
  \AdvePre_\TW(\mathcal{A})
  \;\leq\;
  \Lift_c(\Ntot) + \RootPre_{\tauroot}(\Nprim)
  \;=\;
  \Lift_c(\Ntot) + \frac{\Nprim + 1}{2^{\tauroot}}.
\]
\end{theorem}

\begin{proof}
By the Lift Application corollary (\Cref{cor:lift-application}) for the
preimage game---whose target $T^*$ is fixed before the primitive
sampling---the derived random oracle adversary $\mathcal{B}$ has
$\qRO(\mathcal{B}) \leq \Nprim$, and the random oracle bound of
\Cref{thm:epre-ro} is nondecreasing in $\qRO(\mathcal{B})$, so the monotone
form of \Cref{cor:lift-application} gives the stated bound. With one fixed
target tag, \Cref{thm:epre-ro} applies a candidate-preimage bound against
that target and drops the exponent from $\tauroot(s-1)$ to $\tauroot$.
\end{proof}

\begin{corollary}[Root Tag as a Compact Commitment]
\label{cor:tag-commitment}
The wrapper $\HTW$ is a secure hash of the full encryption context
$(K, U, A, M)$ in the sense of \cite{RS04}: collision-resistant
(\Cref{thm:coll}), preimage-resistant on each fixed input slice
(\Cref{thm:pre}), second-preimage-resistant
(\Cref{cor:second-preimage-rs04}), and everywhere preimage-resistant
(\Cref{thm:epre}). Under the implemented
parameters and work cap $\Ntot \leq 2^{63}$, these advantages are at most
$2^{-128}$ (\Cref{cor:concrete-security}), so the single $\tauroot$-bit
root tag is a binding compact commitment to $(K, U, A, M)$ with standard and
everywhere preimage resistance, native to the mode.
\end{corollary}

\begin{proof}
Collision resistance is \Cref{thm:coll}; standard preimage resistance on
fixed input slices is \Cref{thm:pre}; second-preimage resistance follows from
collision resistance by \Cref{cor:second-preimage-rs04}; and everywhere
preimage resistance is \Cref{thm:epre}. These guarantees quantify the
adversary's success at the same $\tauroot$-bit projection of the final root
transcript, and \Cref{cor:concrete-security} evaluates them at the
implemented parameters. Binding is the compact-commitment property supplied
by collision and second-preimage resistance. Standard preimage resistance is
the property that a tag produced from a uniformly sampled fixed-length
opening does not reveal an efficiently findable opening; everywhere preimage
resistance is the corresponding fixed-target property for tags chosen
independently of the primitive.
\end{proof}

\subsection{Committing Consequences}

\begin{theorem}[Random Oracle CMTDs-4]
\label{thm:cmtds4-ro}
For every $s \geq 2$ and every $s$-way CMTDs-4
adversary $\mathcal{B}$ in the \TW{} random oracle model
(\Cref{sec:ro-model}),
\[
  \Advcmtds{4}_\RO(\mathcal{B})
  \;\leq\; \RootColl_{\tauroot}(\qRO(\mathcal{B}), s)
  = \frac{\binom{\qRO(\mathcal{B}) + s}{s}}{2^{\tauroot (s - 1)}}.
\]
\end{theorem}

\begin{proof}
Run the random oracle CMTDs-4 game. After rejecting on any failed syntax,
message-length, or pairwise-full-input check, the challenger performs $s$
deterministic decryption checks $\Dec_{K_i}^\RO(U_i,A_i,C \concat T)$ on
the one common body $C \concat T$, reaching $R_1,\dots,R_s$. On a winning
execution the tuples $(K_i,U_i,A_i,M_i)$ are pairwise distinct, and since
all checks decrypt the common body, determinism of $\Dec$ makes the
triples $(K_i,U_i,A_i)$ pairwise distinct---equal triples would return
equal messages, hence equal tuples. By Opening Triple Separation
(\Cref{cor:triple-sep}) the candidates $\flatmap(R_1),\dots,\flatmap(R_s)$
are pairwise distinct (leaf tag collisions cannot merge them), and every
accepting check carries the common root tag $T$. The Root Tag Hash and Target-Hit Bound
(\Cref{lem:root-tag-hash}, part~1) gives
$\Advcmtds{4}_\RO(\mathcal{B}) \leq \RootColl_{\tauroot}(\qRO(\mathcal{B}), s)$.
\end{proof}

\begin{corollary}[CMTDs-4 Committing Security]
\label{cor:cmtds4}
For every $s \geq 2$ and every $s$-way CMTDs-4
adversary $\mathcal{A}$ against \TW{} with the resources of
\Cref{sec:resources},
\[
  \Advcmtds{4}_\TW(\mathcal{A})
  \;\leq\;
  \Lift_c(\Ntot) + \RootColl_{\tauroot}(\Nprim, s)
  \;=\;
  \Lift_c(\Ntot) + \frac{\binom{\Nprim + s}{s}}{2^{\tauroot (s - 1)}}.
\]
\end{corollary}

\begin{proof}
A CMTDs-4 winner gives one ciphertext $C \concat T$ and $s$ distinct inputs
$(K_i,U_i,A_i,M_i)$ such that
$\Dec_{K_i}(U_i,A_i,C \concat T) = M_i$ for every $i$. By \TW{} tidiness
(\Cref{lem:tw-tidiness}), each opened input re-encrypts to the same ciphertext:
$\Enc_{K_i}(U_i,A_i,M_i) = C \concat T$. In particular the opened inputs share
the root tag $T$. For the displayed no-$\LeafColl$ bound, use the dedicated
root-collision derivation of \Cref{thm:cmtds4-ro}, lifted by
\Cref{cor:lift-application}: in this all-bodies-equal case, the single shared
body $C$ forces the triples $(K_i,U_i,A_i)$ pairwise distinct; by Opening
Triple Separation (\Cref{cor:triple-sep}), a leaf tag collision cannot merge
two openings and no $\LeafColl$ term arises.
\end{proof}

\paragraph{Remaining committing notions.}
The hierarchy and $s$-way extension of the \cite{BH22} committing notions
are given in \cite[\S3, Figs.~4--5, and App.~A]{BH22}. We use those
relations for the D-notions and a direct source-game argument for the
CMTs-$\ell$ notions to preserve the resource split of
\Cref{sec:resources}.

\begin{corollary}[Committing-Notion Root-Collision Bound]
\label{cor:committing-combined}
For every $s \geq 2$, every
\[
  X \in \bigl\{\mathsf{CMTDs\text{-}1},\, \mathsf{CMTDs\text{-}3},\,
              \mathsf{CMTDs\text{-}4},\, \mathsf{CMTs\text{-}1},\,
              \mathsf{CMTs\text{-}3},\, \mathsf{CMTs\text{-}4}\bigr\},
\]
and every $s$-way $X$-adversary $\mathcal{A}$ against \TW{} with the
resources of \Cref{sec:resources},
\[
  \Adv^X_\TW(\mathcal{A})
  \;\leq\;
  \Lift_c(\Ntot) + \RootColl_{\tauroot}(\Nprim, s)
  \;=\;
  \Lift_c(\Ntot) + \frac{\binom{\Nprim + s}{s}}{2^{\tauroot(s-1)}}.
\]
\end{corollary}

\begin{proof}
CMTDs-4 is bounded directly by \Cref{cor:cmtds4}. For the remaining
D-notions, the $s$-way hierarchy of \cite[\S3, Fig.~5 and
App.~A]{BH22} gives, for $\ell \in \{1,3\}$,
$\Advcmtds{\ell}_\TW(\mathcal{A}) \leq
\Advcmtds{4}_\TW(\mathcal{A})$ for the same output distribution, so the
resource counts are unchanged.

For CMTs-$\ell$, the challenger performs $s$ deterministic encryption
checks, which supply the candidate final-root transcripts
(\Cref{lem:root-candidate-seq} builds the candidate sequence for
encryption and decryption checks alike). On a winning execution all
$s$ encryptions are equal, hence of equal length and with one common root
tag $T$. Pairwise $\WiC_1$-distinctness gives pairwise distinct keys,
pairwise $\WiC_3$-distinctness gives pairwise distinct $(K,U,A)$ triples,
and pairwise $\WiC_4$-distinctness gives pairwise distinct full inputs; in
the last case, equal triples and equal ciphertexts would decrypt
deterministically to equal messages, contradicting full-input
distinctness. Thus each CMTs-$\ell$ winner reaches pairwise distinct
opening triples. Opening Triple Separation (\Cref{cor:triple-sep}) makes
the $s$ final-root candidates distinct while their $\tauroot$-bit
projections all equal $T$.
The hierarchy argument does not change the output distribution or the
resource counts, and the CMTs-$\ell$ encryption checks are the
construction-internal work counted in $N$ by \Cref{sec:resources}.

By the Lift Application corollary (\Cref{cor:lift-application}) for the
$X$ game, the derived random oracle adversary $\mathcal{B}$ has
$\qRO(\mathcal{B}) \leq \Nprim$. The Root Tag Hash and Target-Hit Bound
(\Cref{lem:root-tag-hash}, part~1) then bounds its win by
$\RootColl_{\tauroot}(\qRO(\mathcal{B}), s)$ in the random oracle model,
which the lift transfers to the public permutation setting at $\Nprim$;
the encryption checks are construction-internal work counted in $N$, not
in $\Nprim$.
\end{proof}
